\chapter{1992 Nobel Prize in Economics}
\textbf{Laureates: }Gary Becker (USA)

\section*{Reason for Award}
For having extended the domain of microeconomic analysis to a wide range of
human behavior and interaction, including nonmarket behavior

\section{What They Did}
Gary Becker applied microeconomic analysis to areas previously studied mainly
by sociology, demography, and criminology. His economic approach modeled
individuals and households as purposeful decision-makers choosing subject to
constraints, and tested how well that framework could organize explanations of
diverse social phenomena. Becker's work greatly expanded the boundaries of
economic inquiry and influenced several disciplines.

Becker formalized the economics of human capital, treating education and
training as investments with costs and future returns. This provided a
framework for studying wage differentials, educational decisions, on-the-job
training, and workforce development. Human-capital accumulation became an
important explanation for individual earnings patterns and one factor in
research on national growth differences.

Becker developed an economic model of crime in which a potential offender
weighs expected benefits against costs such as the probability and severity of
punishment. The framework showed how crime may respond to incentives and
influenced subsequent research in criminology and criminal justice.

Becker extended labor and household economics through a theory of household
production, analyzing the allocation of time among market work, leisure, and
domestic activity. He analyzed fertility decisions using opportunity costs and
the quantity--quality tradeoff in investments in children. With Kevin Murphy,
he developed a rational-addiction model in which current consumption depends
on habits and expectations about future costs and benefits.

Becker's analysis of discrimination modeled prejudice as a cost or preference
that can affect market choices. In competitive settings, such preferences can
impose costs on those who discriminate, but discrimination need not disappear
when markets are imperfect or when institutions and information sustain it.
The framework contributed to economic analysis of discrimination; it did not by
itself provide the legal foundation for anti-discrimination or affirmative-action
policies.

\subsection*{Human Capital Theory}
Becker analyzed education, training, and health as investments that can yield
future returns through higher earnings and improved well-being. The human-
capital framework helps explain why educated workers may earn more, why returns
to education vary across fields and people, and how individuals decide how much
to invest in their development. Human capital is now a standard concept in
labor and development economics, although observed earnings also reflect
ability, family background, institutions, and discrimination.

\subsection*{The Economics of Discrimination}
Becker showed that taste-based discrimination can impose costs on discriminators
in competitive markets. Employers who discriminate may pay more for labor or
forgo productive matches. Competitive pressure can reduce such discrimination
under some conditions, but market power, incomplete information, segregation,
and institutional barriers can allow it to persist.

\subsection*{Rational Addiction}
Becker's rational-addiction model treats addictive behavior as a dynamic choice
in which current consumption affects future utility and future consumption. It
offers one way to analyze persistence and quitting decisions, alongside medical,
psychological, and social explanations. It should not be read as a claim that
addiction is simply voluntary or that one policy response is always more
effective than another.

\section{Impact on Economic Thinking}
Becker greatly expanded the scope of economic analysis by applying a purposeful-
choice framework to diverse social phenomena. His work helped establish or
advance the economics of crime, family economics, human capital, and labor-market
discrimination; it was not itself behavioral economics, which often relaxes or
tests the rational-choice assumptions he used.

Crime economics influenced research on sentencing, policing, and correctional
policy. The Becker framework predicts that deterrence can affect crime by
changing expected costs, but actual effects depend on enforcement, behavior,
social conditions, and the design of punishment.

Family economics informed research on household behavior and welfare policy.
Becker's fertility analysis highlighted opportunity costs and the
quantity--quality tradeoff in investment in children, contributing to research
on demographic change. Policy outcomes also depend on culture, institutions,
and social norms.

Human-capital frameworks became standard in analyzing education, training, and
workforce development. They organize research on returns to schooling, wage
inequality, migration, health, and economic growth, while recognizing that
education is not the only source of observed earnings differences.

\subsection*{Becker's Methodological Legacy}
Becker's interdisciplinary modeling approach encouraged economists to engage with
complex social questions beyond traditional market boundaries. He showed that
economic methodology is a powerful tool for understanding a broad spectrum of human
behavior. His work bridged economics, sociology, and psychology, creating richer
analytical frameworks for studying human social interaction.

\subsection*{Expanding the Economic Horizon}
Becker demonstrated that economics is not limited to market transactions but
provides tools for analyzing any situation involving choice under constraints.
This expansion of economic methodology legitimized the study of non-market
behaviors and created new research frontiers. Subsequent economists have extended
Becker's approach to marriage, divorce, household dynamics, and social networks.

\section{Modern Perspectives Today}
Behavioral economics tests and extends the rational-choice framework associated
with Becker by incorporating bounded rationality, preferences, and other
psychological factors. Crime economics continues to inform research on
deterrence, incapacitation, and rehabilitation, while recognizing limits to
simple incentive models.

Research on family formation and household decisions uses opportunity costs,
time allocation, and the economics of childbearing developed in this tradition.
Human-capital calculations inform some analyses of education and training
policy, while discrimination enforcement also relies on legal, institutional,
and social evidence beyond economic models.

Research on addiction sometimes uses dynamic-choice models to study dependence
and habit formation, alongside clinical and social-science approaches. Related
economic methods are used to study migration, political behavior, and human-
capital investment in development, without implying that all such behavior is
fully rational or that Becker originated each application.

Gender economics uses household-production and time-allocation models to study
labor-force participation, wage differentials, and the division of household
labor. Research on the care economy likewise recognizes unpaid work as
economically significant, while demographic models extend fertility analysis by
incorporating cultural factors and institutional constraints.

Becker's methodological legacy is evident in the pervasive application of
economic reasoning to social phenomena. His expansion of economic methodology
continues to inspire researchers probing deeper connections between individual
choices and social outcomes. The Becker tradition of applying rigorous economic
analysis to seemingly non-economic domains remains one of the most influential
approaches in modern social science.

\subsection*{Key References}
\begin{itemize}
    \item Becker, G. S. (1964). \textit{Human Capital: A Theoretical and Empirical Analysis, with Special Reference to Education}. Columbia University Press for the National Bureau of Economic Research.
    \item Becker, G. S. (1981). \textit{A Treatise on the Family}. Harvard University Press.
\end{itemize}
